% !TEX root = Hogwarts_1347.tex

\chapter{The Tide of the Great Hulks}

They exited Diagon Alley through the same unassuming door through which they had entered — wedged between two shops and marked only by the small, carved symbol of the Crossed Keys. Thomas pulled it open, and they found themselves back in the long, warped corridor of the bridge's interior. Navigating past the myriad of doors, they reached the portal painted with the same keys. As they crossed back into the cramped airlock, the familiar scents of stale ale and freshly baked bread rushed to meet them. The space was now agonisingly tight, with the large leather trunk and Pippin's cage barely squeezing in alongside them. Emerging onto the London Bridge, the bell of the Chapel of St. Thomas rang out four sharp peals, echoing over the Thames. Given the late hour and their heavy load, Alianor hailed a muggle carriage. They rattled northward across the bridge before turning eastward, following the river's edge toward the docks.

The carriage deposited them near the bustling Billingsgate market. They navigated a maze of back alleys, turning left and right past the smell of salt and scales until they reached a peculiar wall---a patchwork of timber-framed wattle-and-daub and old, repurposed, Roman red bricks. Stepping through the solid surface, they emerged onto the wizarding margins of the Thames.

Five massive vessels were moored to the stone pier. The first four, identical sister ships named the \textit{Mermaid}, the \textit{Siren}, the \textit{Ceasg} and the \textit{Melusine}, were behemoths of naval architecture --- a magical hybrid of the sturdy cog and the spacious hulk, but nearly twice the size of any merchant vessel Thomas had ever seen. Their hulls were deep and rounded, featuring high, fortified forecastles and even more impressive sterncastles that towered like wooden fortresses over the water. Each sported two masts with immense square sails.

The fifth ship, the \textit{Merrow}, was a relic of an older age. It was a long, low-slung galley, roughly forty meters in length, with thirty-two heavy oars protruding from each side like the legs of a centipede. It lacked the high castles of the sister ships, offering smaller living quarters, its two sails striped in the colours of all four Hogwarts houses.

Near the pier stood a small stone-and-wood office where a goblin clerk sat behind a high desk. He ignored the frantic activity of house-elves loading barrels and pallets of supplies, his quill scratching rhythmically as he processed passengers. When they reached the front, Thomas placed a gold coin on the counter.

\enquote{Single ticket for Hogwarts,} Thomas said.

The goblin took the coin without a hint of emotion. \enquote{First year?} he asked rhetorically. Thomas nodded. \enquote{Better to avoid the Slytherin lot then,} the goblin grunted, sliding back several silver coins and a handwritten slip of paper. \enquote{You're on the \textit{Siren}, middle deck, cabin eleven.} He gestured toward the second ship. At the prow, carved directly into the stempost and projecting along the beak beneath the forecastle, a siren of dark oak leaned forward with the rake of the bow, her upper body human and her scaled tail drawn tight against the timber, cut in shallow relief so the sea could wash over her without tearing the carving away. \enquote{The ship convoy will stop at Dublin, the Isle of Man, and Dumbarton - on the Clyde,  some miles downstream from Glasgow} the goblin explained, his voice a dry rasp. \enquote{Then it will head north, hugging the west coast of Scotland. From there, you'll travel by carriage to Hogsmeade and finally Hogwarts castle.} He tapped the parchment. \enquote{You'll have to share the berths section with only one other student for now; two others are scheduled to embark at the Isle of Man and Dumbarton.}

\section*{The Siren's Call}

Thomas took the paper, his heart hammering against his ribs. At the foot of the \textit{Siren's} gangplank, Alianor stopped. The iron mask of her composure finally slipped, her eyes glassing over with unshed tears. She pressed the leather purse of coins into his hand.

\enquote{For the rest of your uniforms once you are sorted, and for your time in Hogsmeade,} she whispered, her voice thick with emotion. She reached out, straightening his collar one last time. \enquote{Be presentable, Thomas. For your professors and your classmates. Eat well, study hard, and\ldots{}} she looked at the owl in its cage, \enquote{\ldots{}write to me. Often.}

Thomas reached for his heavy trunk and Pippin's cage, preparing for the arduous climb, but a house-elf nearby intervened with a quick flick of its hand. A subtle shimmer of magic washed over the luggage, and both the trunk and the cage began to levitate, bobbing weightlessly in the air. With significantly less effort, Thomas simply guided them through the narrow gangway. He watched the trunk drift ahead of him, thinking to himself that such a spell would be a convenient one to master.

\section*{Cabin Eleven}

Stepping off the gangway and onto the \textit{Siren's} open wooden deck, Thomas found himself in a whirlwind of activity. Two wizards stood mid-ship, their wands tracing arcs through the air as they directed barrels and pallets of supplies toward the lower decks with precision. The crowd of travellers split; some headed toward the towering forecastle at the bow, while others moved toward the sterncastle. One of the wizards, holding his wand pressed firmly against his own throat to amplify his voice, barked orders over the din, directing passengers to their respective cabins located in the castles and the more cramped quarters deep within the hull.

The hierarchy of the ship was reflected in its architecture. The sterncastle housed the most prestigious quarters---wood-panelled, private cabins for Hogwarts staff and professors, distinguished guests, and the families of the school's patrons, complete with stern-walks overlooking the ship's wake. Below them, and forward in the forecastle, senior students shared four-berth cabins with solid doors and the luxury of portholes.

Thomas, however, was directed toward the belly of the ship. The younger students were housed in cavernous spaces on the middle and lower decks, just above the cargo hold. Here, the layout consisted of long corridors with walls that separated the open cabins on each side. Each cabin was remarkably narrow, containing four berths---two stacked atop two---perpendicular to the corridor. A tiny, foldable wooden table, held in place by a simple cord, was the only furniture other than the beds. There was barely enough room under the lower berths for four trunks, and with a ceiling only two meters high, the taller students were forced to bow to avoid the heavy, exposed support beams.

Despite the lack of portholes, the middle deck felt surprisingly welcoming. It was decorated in a mixture of polished, varnished wood and clean white panels with brass accents. Oil lamps bathed the space in a warm glow, and gaps left above and below the curtains allowed for essential ventilation. In the common area of the main deck below the sterncastle, passengers gathered at fixed tables to share meals or pass the time with backgammon, ludus latrunculorum, pachisi or wizard's chess.

Though it was a massive vessel, the \textit{Siren} required a crew of only eight to twelve wizards --- far fewer than a Muggle ship of its size. Spells significantly improved ship performance by offering higher, more consistent wind speeds, resulting in much greater vessel velocity and shorter journey durations. Additionally, magic provided vital concealment, hiding the convoy from Muggle ships and ports.

Thomas navigated a steep, narrow flight of stairs, thankful once again for the levitation spell on his trunk. He reached a narrow corridor illuminated by small hatches that let in the fading afternoon light, supplemented by rows of oil lamps. The air here was pleasant, carrying the fragrance of olive oil scented with lavender. He eventually found his quarters, marked with an engraved wooden sign bearing the number 11.

Pulling back the thick curtain, he found the room lit by two oil lamps and four beeswax candles in brass candelabras. On the bottom right berth, he immediately noticed a broom, similar to the one he saw earlier at Diagon Alley, wrapped with a red knot and a letter addressed to him---his \textit{Wayfarer}.

One of the cabin's occupants was already there. A boy, perhaps an inch taller than Thomas with short black hair, hopped down from the top left berth and extended a hand.

\enquote{Baldwin\ldots{} Baldwin Attewell. Nice to meet you,} the boy said with a friendly grin.

\enquote{Thomas Blackwood, nice to meet you,} Thomas replied.

Baldwin's eyes widened, and his voice rose in surprise. \enquote{Blackwood? As in Mary Elizabeth Blackwood? The Auror of the Council and the youngest Hogsmeade tourney champion ever?}

Thomas offered a faint smile and a nod. \enquote{She's my aunt, my father's younger sister.}

\enquote{If you don't mind me saying,} Baldwin remarked, looking him over, \enquote{you don't exactly resemble her,} clearly referring to Mary's legendary long, white-blonde hair and her strikingly pale skin marked with white freckles.

In response, Thomas raised his left hand and ran his fingers through his short, dark brown hair. As he did, the strands shimmered and lengthened, turning a bright, platinum blonde. His skin paled instantly, blooming with the delicate white freckles that were the hallmark of the Blackwood family. He paused for a moment, then brushed his hand through his hair again; the colour shifted back to dark brown, and the sickly pallor of his skin warmed back to a healthy, almost tanned glow.

\enquote{Wow,} Baldwin breathed, leaning in. \enquote{You must teach me how to do that.}

\enquote{I would love to,} Thomas confessed, sliding his trunk under the berth and hanging Pippin's cage on one of the wall hooks, \enquote{but I don't actually know how. I just\ldots{} do it.}

As he settled in, Thomas noticed a white, fluffy shape curled up under the opposite berth. Baldwin followed his gaze.

\enquote{Don't worry about Dolores,} Baldwin said, nodding toward the white cat. \enquote{She's far too slothful to even bother attacking him\ldots{} or is it her?}

\enquote{He's a him,} Thomas confirmed with a smile, watching the young tawny owl adjust to the swaying of the ship. \enquote{I call him Pippin.}

The boarding process was nearly complete, and as students settled into their respective quarters, the corridors below decks grew visibly quieter. Across the narrow passageway from Cabin 11, two young witches arrived at their own doorway, struggling slightly with two heavy trunks---one a cheerful yellow and black, the other a shiny, deep brown that looked brand new. Thomas immediately recognised them as the sisters from the tailor shop, Cecile and Celine.

In that moment, Pippin, perhaps spurred by the excitement of the journey, managed to nudge the latch of his cage loose. With a sudden flutter of soft feathers, he took flight, swooping across the corridor and landing squarely on the shoulder of the younger sister, Celine. She had been mid-bite into a small tourtelette, and the little owl was staring at it with unblinking intensity.

Celine laughed, her French accent lilting through the air, and pinched off a piece of the golden crust. Pippin devoured it instantly. She then reached up to stroke the top of his head; the young tawny owl stretched his wings and leaned into her touch, seemingly melting against her shoulder as he enjoyed the affection as much as the snack.

Thomas reached the doorway of his cabin, watching the scene with a mix of embarrassment and amusement.

\enquote{I see he has a fine taste,} Celine joked, looking at Thomas. \enquote{Apparently, even the owls of London recognise a proper Lyonnais tourtelette.}

Cecile smirked, nudging her sister. \enquote{You and your new feathered friend are due for a disappointing surprise once you see the ship's mess menu a few days from port.}

Thomas leaned against the doorframe, shaking his head. \enquote{Wow, he's been bought off already. Absolutely shameless,} he joked. \enquote{Just a piece of pie and he goes all soft. Look at him, he's melting already.}

\enquote{Don't forget the head-scratches,} Celine reminded him, her fingers still buried in Pippin's soft plumage.

Thomas looked the owl in the eye and laughed. \enquote{Traitor. So much for loyalty.}

\enquote{I remember seeing you at the tailor's in the Alley,} Thomas added as the girls finally shoved their trunks into their narrow cabin.

\enquote{And we remember you with your mother,} Cecile replied, smoothing out her robes.

\enquote{Would you like to join us in our cabin for a bit?} Thomas asked, gesturing toward Baldwin. He looked at Cecile, the elder sister. \enquote{It would be nice to hear from someone who actually knows what to expect at Hogwarts.}

After securing their luggage under the lower berths, the sisters accepted the invitation and followed Thomas across the hall. Inside Cabin 11, Baldwin stood up, offering a polite bow as they introduced themselves.

\enquote{Cecile Fournier,} Cecile said, introducing herself, before adding and gesturing to her sister, \enquote{and this is Celine.}

\enquote{Baldwin Attewell,} Baldwin replied, catching the cadence of their speech. \enquote{That's a fine accent. Where are you from originally?}

\enquote{Lyon,} Celine answered proudly. \enquote{It sits where the Rhône and the Saône rivers meet, far to the south and east.}

Baldwin nodded, showing a surprising grasp of geography. \enquote{I know the place. Famous for its wine and its hills. But that raises a question---why come all the way to the North? Why Hogwarts and not Beauxbatons?}

Cecile's expression turned a bit more serious, her gaze lingering on the flickering oil lamp. \enquote{That,} she said softly, \enquote{is a long story.}

\section*{The Merchant and the Baker}

\textit{Lyon, Five Years Prior}

The sun hung low over the hills of Fourvière, casting long, amber shadows across the terracotta roofs of Lyon. Janus Laskaris walked with a steady, rhythmic gait, his young daughter Philippa skipping at his side. They began their journey on the rue Mercière, the heart of the city's commerce, where the air was thick with the scent of tanned leather and expensive spices. Turning toward the river, they crossed the Pont Change, a massive stone bridge lined with the shops of money-changers and goldsmiths, their scales clinking in a metallic chorus above the rushing Saône. As they reached the western bank, the streets of Vieux Lyon narrowed into a labyrinth of medieval stone. The buildings here were tall and proud, featuring the famous \textit{traboules}---secret passageways that cut through courtyards and under spiralling stone staircases. They walked toward the Montée des Carmes Déchaussés, where the incline of the hill began its steep ascent.

Janus Laskaris was a man of the world in the most literal sense. A fabulously rich merchant of both common and magical cargo, he was no stranger to the great centres of civilisation. He maintained residences in London, Paris, Rome, and Constantinople, and his business frequently took him to the markets of Venice, Athens, Alexandria, Jerusalem, Damascus, and Isfahan. Yet, despite his global reach, Lyon held a unique resonance for him. Although it was a profitable commerce hub, his reason for choosing it as a permanent residence was deeper; something about the city simply screamed \enquote{home.} The food and wine were incomparable, and the climate suited him perfectly. Even as a non-native speaker, he found that the Lyonnais French carried a rounded warmth; the words flowed more gently than the sharper speech of Paris, with vowels that lingered and consonants that softened. He took great pleasure in both speaking and listening to the local tongue.

Janus was accustomed to the company of the powerful, counting monarchs, ministers, doges, and emperors among his acquaintances. In the hidden world, he was well-known to the members of various wizarding councils, synods, and majlis, and he took pride in his contacts among eminent professors and headmasters. Yet, among all these prestigious figures, he particularly enjoyed the company of a humble Muggle resident: Monsieur Henry Fournier. Henry was a baker whose shop was located on the west margin of the Saône, near the stairs of the Montée des Carmes Déchaussés. While Janus could easily have had the finest pastries delivered to his mansion or prepared by his own staff, he preferred to visit the bakery in person. He enjoyed hearing the local news and gossip --- who was marrying whom, who was expecting a child, whether the tavern owner had finally been caught diluting his wine, or how many men had been requisitioned for King Philip's army. In return, the young baker was fascinated by Janus's tales of exotic lands like Rome, Constantinople, Jerusalem, and Isfahan, captivated by the different customs and extravagances of distant rulers.

However, the nature of their friendship was about to change. On this particular day, at his daughter's request, Janus had brought high-quality ingredients so the baker could make a delicacy from his home: Michaelmas bread loaves. As the merchant and the baker sat sharing a loaf and some Rhône valley wine, discussing the latest local rumours, Janus noticed something peculiar. He was seated at a table facing the opening of the oven, while Henry sat opposite him with his back to the heat. Closer to the oven was Cecile, the baker's eldest daughter, who was about eight years old. Janus watched in silence as the glowing coals inside the oven began to coalesce; the embers rose and swirled, forming the distinct shape of a bird in flight. It flapped its fiery wings for a heartbeat before dissolving into a blooming flower and then a six-pointed star.

Janus glanced at his own daughter, wondering if she were the source of the magic. Philippa was on the floor with the baker's youngest daughter, Celine, who was two years younger than Cecile. They were drawing on a sheet of paper he brought from the orient with paints from the Italian peninsula. Janus felt a brief sense of relief when he saw Philippa drawing the same shapes on the paper as those appearing in the oven, assuming it was her influence. However, his relief was short-lived. He noticed a rash on Cecile's forearm and asked her about it. When the girl mentioned she had scratched it on the leaves of a rowan tree, Janus became instantly alert. While harmless to Muggles, the untreated wood and leaves of the rowan are a severe irritant to wizards. Though its treated wood is prized for wands --- especially for defence against the dark arts --- the rowan's fruits, flowers, and sap are dangerous and potentially fatal to wizards.

Later that night, Janus sent an owl to the \textit{Confrérie des Arts de la Sorcellerie de France}, the prestigious wizarding guild in Paris. His contact there consulted \textit{Le Livre d'Admittance}, the book recording every young wizard in the region defined by the six corners between Bretagne, Flandre, Alsace, Provence, Roussillon et Guyenne. Two nights later, the answer arrived, and Janus was stunned. As expected, Cecile Fournier's name was in the book, but to his immense surprise, so was the name of Celine Fournier.

Janus knew that the world was about to be turned upside down for the Fourniers. Uncomfortable as it was, he felt it was his duty to present the wizarding world to them himself, rather than leaving it to a clerk from the \textit{confrérie} or a faceless letter. He realised he would miss the simple talks about Muggle life he shared with Henry; their conversations were about to become much more serious.

Janus took his wand, a small locket, and the letter from the Confrérie and headed to the bakery. Henry was surprised to see him so early, and looking so somber. Janus read the letter aloud to the trio, who were left entirely confused.

That confusion turned to sheer terror when Janus drew his wand. With a sharp gesture, Janus ignited the cold charcoal in the oven; with a second, he levitated a raw loaf of bread through the air and tucked it neatly onto the baking stone. Henry instinctively scrambled backward into a corner, pinning himself against the wall and shielding his daughters behind him. With a calm and serene tone, Janus assured the baker that he was still the same person Henry had known for years --- the man who had introduced him to his late wife and mother of his daughters and later organised her funeral.

Intrigued, Celine approached the table where Janus had placed an open locket. Inside was a portrait that she could swear had just moved; the figure in the painting adjusted herself and then walked out of the frame entirely. Janus explained that this was a wizarding painting. He then assured Mr. Fournier that any rumours he might have heard in sermons or old tales --- claiming that witchcraft was an unnatural or inherently evil art --- were completely unwarranted. This seemed to calm the baker, who eventually sat back down at the table with Janus, his oldest daughter beside him and the youngest on his lap.

Janus then proceeded to reveal the true scale of the wizarding world. Using his wand to command charcoal dust and flour to form images and shapes in the air, he showed a map of the known world and its most important wizarding settlements. He pointed out eight eminent schools marked by incandescent charcoal:

\begin{itemize}
    \item {Hogwarts}: in the Highlands of Scotland.

    \item {Beauxbatons}: within the French Pyrenees.

    \item {Durmstrang}: in the north of Scandinavia.

    \item {Uagadou}: in Buganda.

    \item {Mahoutokoro}: in the far East.

    \item {Koldovstoretz}: along the banks of the Upper Volga.

    \item {Kamryn}: in the Carpathian Mountains.

    \item {Malakopí}: within the massive underground community in Cappadocia.
\end{itemize}

Janus explained that in most of these schools, young wizards studied from ages eleven to seventeen. He offered to provide everything they would need regardless of their choice, but admitted he had more influence at Hogwarts, where he had studied and where he intended to send Philippa. As a patron of the school and member of its board he had helped recruit professors and staff from across the globe. He noted that while Beauxbatons was perhaps stronger in potions and alchemy due to a rising professor named Nicolas Flamel --- whom he regrettably had failed to recruit --- Hogwarts possessed the finest master of magical artefacts in Professor Magnus Grimm, whom Janus had brought from Constantinople. He also promised to provide tutors to ensure the girls were literate in Latin or Greek before they attended.

Henry hesitated, not out of fear of magic, but because of the vast distance. He also questions the travel and the difficulties given the conflict between England and France. Janus remained somewhat vague but stated that relations in the wizarding world were currently less conflicting. He gave his personal guarantee that he would personally bring the girls back to Lyon every summer in his own ship---a vessel that was part of a fleet protected from Muggle eyes and capable of navigating both the high seas and the Rhône river.